\chapter[永远，在某处，有人在为你而战]{永远，在某处，有人在为你而战}

\chaptersubtitle{《魔法少女织莉子》}

与《魔法少女和美》一样，漫画《魔法少女织莉子》（Mura Kuroe 作画，原案出自 Magica Quartet；这一集体署名指的是《魔法少女小圆》的制作人岩上敦宏、导演新房昭之、编剧虚渊玄以及角色设计苍树梅）也必须面对由本系列三条无法回避的真理所构成的陷阱。
简要回顾一下：
魔女既是对支配性权力结构的挑战，又是一种腐败因素；
以丘比为代表的这一支配性权力结构极其不公；
而在外传作品中，这一权力结构又不可能被推翻，因为这种推翻发生在《魔法少女小圆》最终话之中。
《魔法少女和美》强调的是第一条原则，因此落入了将反抗体制者污名化的陷阱。
《魔法少女织莉子》则强调第三条原则，因此避开了这一陷阱，但代价是自身成为一出悲剧。

与《魔法少女小圆》类似，《魔法少女织莉子》也分为三个叙事段落，并由标题人物织莉子将其串联起来。
第一段在故事内外都是一枚误导性的烟幕弹：
到了第三段才揭示出来，织莉子为了转移丘比的注意力，操纵了一个遭受虐待、又失去双亲的幼童由麻，使她与丘比签订契约，成为魔法少女。
对读者而言，这会制造出一种预期：
这部漫画大体上会像前两话实际呈现的那样，讲述杏子与由麻之间的友谊以及导师—弟子的关系。
然而，第二段叙事却推翻了这一预期。
这个故事在第二话结尾戛然而止；
第三话和第四话转而跟随麻美，讲述她调查魔法少女连续遇害事件，并与凶手，也就是织莉子的朋友、同为魔法少女的雾香交战。
尽管雾香拥有减缓时间的能力，麻美最终仍然将她重创；
雾香随后逃回织莉子身边。
最终段落发生在见泷原中学。
织莉子和雾香占据学校，并使整所学校都笼罩在魔女结界之中。
这一段主要跟随焰展开：
我们得知，织莉子的真正目标是通过此刻杀死小圆，拯救世界免于毁灭，因为小圆日后终将变成毁灭世界的魔女；
而焰当然很快便与织莉子对峙，并同她战斗。
雾香变成了魔女。
漫画没有明说，但暗示她在与麻美一战之后就已经开始转化，而笼罩学校的结界以及袭击学生的使魔都属于她。
麻美、杏子和由麻则与她交战。
最终，焰成功杀死织莉子；
但织莉子在最后时刻杀死了小圆，迫使焰重置时间线。

这一故事中处处都是悲剧。
几乎每个角色都被看似积极的冲动所毁灭。
最明显的是，焰未能保护小圆，而织莉子走向死亡，原因正在于两人都决心保护对自己而言重要之物。
织莉子无法冒险让小圆毁灭世界，因此在她看来，除了试图杀死小圆之外，没有任何可行的行动路径；
焰无法冒险让小圆受到伤害，因此在她看来，除了消灭织莉子之外，也没有任何可行的行动路径。
结果是两人都失败了：
焰重置了时间线，并最终通向电视动画的时间线、通向小圆对历史的改写，也通向所有曾经存在魔女的时间线的毁灭，其中自然包括织莉子的时间线；
与此同时，织莉子杀死了她所在时间线中的小圆，使这一轮时间线成为焰的又一次失败。
但悲剧同样也降临在其他角色身上。
小圆之所以会来到织莉子能够杀死她的位置，只是因为她想帮助焰；
如果她留在焰为她设置的保护结界之内，她本可以活下来。
杏子试图保护由麻，使她远离魔法少女生活的危险；
但在这一过程中，她反而让由麻担心自己毫无用处。
于是，由麻为了对杏子有用而成为魔法少女。
而雾香原本只是希望自己拥有不再孤身一人的勇气，却在这个过程中变成了捕食者。

但在这一切之中，最具悲剧性的或许是能够预见未来的织莉子。
漫画以一个相当暧昧的场景结束：
雾香走近一个看似心情低落的织莉子，后者透露自己已经杀死了许多人。
雾香提出要分担她的重负，两人随后手牵着手离去。
Puella Magi Wiki 将这一场景解释为另一条时间线，或者死后的世界\cite{puella-magi-wiki2015-oriko}。
不过，考虑到漫画此前不断穿插织莉子与雾香如何相遇、如何成为魔法少女的闪回，更有可能的是，这一场景也是类似的闪回；
或者更准确地说，漫画中那些“现在时”的事件，乃是织莉子刚刚看到的关于自身未来的幻象。
她看到自己将杀死许多人，也看到自己和雾香都会死去；
但她同时也看到，这意味着她所在时间线的地球将免遭小圆魔女形态的毁灭。
于是，她决定这一切是值得的。

如果最终场景呈现的确实是这样的情形，那么织莉子在整个系列中的平静决意便获得了解释：
她只是在一出早已写好的戏剧中扮演自己的角色。
她知道自己会在何时、以何种方式死去，并且早已决定为了世界的利益而接受这一命运。
在这个意义上，她成了一个堕落的小圆，正如雾香是一个堕落的焰。
雾香与焰的共同点在于，她们都彻底献身于自己所爱的人，并愿意为对方做任何事；
但焰通常会避免伤害人类或魔法少女，而是把目标集中在与魔女和丘比的战斗上，
雾香则猎杀魔法少女，并且只有在直接保护织莉子时才会与魔女战斗。
同样，织莉子和小圆都愿意牺牲自己来拯救世界；
但织莉子愿意杀人，而小圆的目标则是拯救所有人。

由此可见，这就是《魔法少女织莉子》对前述陷阱给出的解决方案：
既然魔女与魔法少女的不公体制不可能在这部漫画中被打破，作品便转而聚焦于一次规模较小的拯救世界的努力。
它展示了每个人如何都注定被周遭体制所毁灭，并将她们未能推翻这一体制的失败，描绘为所有相关人物共同卷入的悲剧。
不过，这部漫画也避开了另一个本可能由此展开的陷阱，即暗示这一体制不可打破。
它通过有意呈现焰重置时间线，并以雾香和织莉子共同分担重负这一相对幸福的结尾收束，避免了这种暗示。

焰仍在继续努力，走向我们在电视动画中看到的结局。
在那之前，彼此之间仍然可以找到慰藉。
只要丘比的系统依旧存在，所有魔法少女终有一天都会被它摧毁；
但正如由麻所说，“终有一天，不是今天。”
